\chapter{Markets, Institutions, and Economic Ideologies}
\label{chap:ideologies}

\section{Learning objectives}

After completing this chapter, you should be able to:
\begin{itemize}
  \item explain the economic case for markets;
  \item identify conditions under which intervention may be justified;
  \item distinguish efficiency, liberty, and equality as evaluative criteria; and
  \item use economic reasoning in ideological debates without reducing them to economics alone.
\end{itemize}

\section{The case for markets}

Markets can coordinate vast numbers of decisions through voluntary exchange and prices. Supporters emphasise decentralised knowledge, incentives for innovation, freedom of choice, and the difficulty of central planning. Where property rights are secure and competition is effective, prices can direct resources toward uses that consumers value and away from uses that are costly.

\section{The case for intervention}

Interventionist arguments begin with the limits of this mechanism: market power, externalities, public goods, information asymmetry, incomplete insurance, and inequality can cause outcomes that fail to meet social objectives. Government may regulate, tax, subsidise, provide services, redistribute income, or create institutions that make markets work better.

Government action also has costs. Policymakers may lack information, face political incentives, be captured by special interests, or administer policies poorly. The relevant comparison is not between an ideal market and an imperfect state, but between feasible institutional arrangements.

\section{Efficiency, liberty, and equality}

Economic analysis identifies trade-offs but cannot rank all values by itself. A policy may improve efficiency yet reduce a person's freedom of contract; another may reduce inequality but weaken work incentives. PPE reasoning requires making these value conflicts explicit, examining empirical consequences, and considering the legitimacy of decision-making institutions.

\section*{Chapter Review}

\subsection*{Key concepts}
\begin{itemize}
  \item \textbf{Decentralised knowledge}: information dispersed among individuals and firms.
  \item \textbf{Government failure}: policy outcomes made worse by information, incentive, or political problems.
  \item \textbf{Institutional comparison}: evaluation of realistic alternatives rather than idealised extremes.
\end{itemize}

\section*{Exercises}

\subsection*{Conceptual questions}
\begin{enumerate}[label=\arabic*.]
  \item Give three economic arguments in favour of markets.
  \item Give three market failures that may justify intervention.
  \item Why is government failure relevant to evaluating policy?
\end{enumerate}

\subsection*{Graphical exercises}
\begin{enumerate}[resume,label=\arabic*.]
  \item Use a supply-and-demand diagram to illustrate one policy that may improve efficiency when a market failure is present.
\end{enumerate}

\subsection*{Analytical exercises}
\begin{enumerate}[resume,label=\arabic*.]
  \item Compare a pollution tax, direct regulation, and no policy in terms of information requirements and incentives.
  \item Explain why an efficient allocation may still be rejected on grounds of distributive justice.
\end{enumerate}

\subsection*{Exam-style multiple-choice questions}
\begin{enumerate}[resume,label=\arabic*.]
  \item The appropriate comparison in policy analysis is generally between: (a) an ideal market and an ideal government; (b) feasible institutional alternatives; (c) economics and no evidence; or (d) equality and efficiency only.
  \item Government failure refers to: (a) every policy that changes prices; (b) problems caused by incentives, information, or politics in public action; (c) all taxation; or (d) competitive equilibrium.
\end{enumerate}

\subsection*{PPE discussion question}

\noindent ``The market versus the state'' is a false choice. Discuss this claim, using examples of institutions that enable markets as well as institutions that correct their failures.
